\part{Complete layout examples}

\section{A complete benchmark table}
Use shared scales within each column, explicit units and a stated ranking rule.
This example combines a score, latency, before/after change and recent history.
Best score and lowest latency are bold; missing observations remain missing.
\begin{demo}{Model comparison with a grouped header}
\begin{gradtable}[width=\linewidth,header align=c,
  label width=8mm,column sep=3pt,min text width=12mm,
  min graphic width=15mm,row stretch=1.7]
  \gradcolumn{method}{Method}
  \gradcolumn[type=bar,mark=best,weight=1.2,
    options={max=100,precision=1}]{score}{Score}
  \gradcolumn[type=number,width=15mm,mark=best,
    options={better=lower,precision=1}]{time}{ms}
  \gradcolumn[type=dumbbell,options={max=100,precision=1,
    compare label=delta}]{change}{Change}
  \gradcolumn[type=spark,options={spark range=fixed,
    min=0,max=100,precision=1}]{trend}{Five periods}
  \gradheader{\gradspan{1}{Setup}\gradspan{2}{Performance}
    \gradspan{2}{Comparison}}
  \gradrow{A,82.4,14.8,68/82.4,{68,74,70,78,82.4}}
  \gradrow{B,91.6,18.2,76/91.6,{76,81,78,87,91.6}}
  \gradrow{C,87.3,11.5,80/87.3,{80,72,83,78,87.3}}
  \gradrow{D,78.9,9.6,83/78.9,{83,72,79,74,78.9}}
\end{gradtable}
\end{demo}
The dumbbell connects reference and current values, and its label shows their
difference. All trends share 0--100 vertically and equal observation spacing.
These numbers are demonstrations, not experimental findings.
